\section{Alineamiento ciclotomico, orbifold of order \(3\) and Monstrous Moonshine}
\label{p12-83-c20-s8:chap:alineamiento-orbifold-moonshine}

The Leech lattice constructed in
\cref{p12-83-c20-s5:thm:identificacion-vecino-leech} makes it possible to pass from finite incidence
to an infinite graded structure. This transition will not be effected
through a chain of names. The following will be constructed, in typed order:
\begin{enumerate}
 \item the intertwining of the twelve fibres \(A_2\);
 \item a fixed-point-free lattice isometry of order three;
 \item its graded trace and its type-zero lift;
 \item the cyclic orbifold and its dual automorphism of class \(3B\);
 \item the moonshine vertex operator algebra;
 \item through distinct publications, its automorphism group, its graded
       character, and the McKay--Thompson series.
\end{enumerate}
Moonshine sera the theorem that coordina the action and the series. is not usara
the atajo not tipado
\(V^\natural\to\mathbb M\to\text{Moonshine}\).

\subsection{Alignment of the \(12\) APP and Witt fibers}

The twelve fibers of \(W_{\mathrm{APP}}\cong A_2^{12}\) carry three
independent labels:
\[
 i\in\mathbb Z/3\mathbb Z
 \quad\text{(coordinate pair)},\qquad
 \mu\in\{0,1\}
 \quad\text{(sheet)},\qquad
 \varepsilon\in\{0,1\}
 \quad\text{(cyclotomic orbit)}.
\]
Fixing the dodecaphasic origin constructed on the TPK, we define
\begin{equation}
 \lambda(i,\mu,\varepsilon)
 :=4i+2\mu+\varepsilon\pmod{12}.
 \label{p12-83-c20-s8:eq:alineamiento-doce-fibras-u031}
\end{equation}

\begin{lemma}[Bijectivity of the alignment]
\label{p12-83-c20-s8:lem:biyectividad-alineamiento-u031}
The map \(\lambda\) from
\eqref{p12-83-c20-s8:eq:alineamiento-doce-fibras-u031} is a bijection onto
\(\mathbb Z/12\mathbb Z\).
\end{lemma}

\begin{proof}
For each \(i\), the four values \(4i+2\mu+\varepsilon\) form the block
\(\{4i,4i+1,4i+2,4i+3\}\). The three blocks for \(i=0,1,2\) are disjoint
and cover the twelve classes.
\end{proof}

Let \(C\) be the Coxeter element of \(A_2\) and let \(\mathsf R\) be the reflection that
satisfies
\begin{equation}
 \mathsf R C\mathsf R=C^{-1}.
 \label{p12-83-c20-s8:eq:reflexion-conjuga-coxeter-u031}
\end{equation}
Once the trivializations \(\iota_b\) and the oriented frames \(P_b\) have been fixed,
we set
\begin{align}
 g_b&:=\iota_b^{-1}P_bCP_b^{-1}\iota_b,
 \label{p12-83-c20-s8:eq:coxeter-calibrado-fibra-u031}\\
 T_{i\mu\varepsilon}
 &:=\iota_{\lambda(i,\mu,\varepsilon)}^{-1}
 P_{\lambda(i,\mu,\varepsilon)}\mathsf R^\mu.
 \label{p12-83-c20-s8:eq:entrelazador-fibra-u031}
\end{align}

\begin{theorem}[APP Entanglement--Witt]
\label{p12-83-c20-s8:thm:alineamiento-app-witt-u031}
For each terna \((i,\mu,\varepsilon)\),
\begin{equation}
 g_{\lambda(i,\mu,\varepsilon)}T_{i\mu\varepsilon}
 =T_{i\mu\varepsilon}C^{(-1)^\mu}.
 \label{p12-83-c20-s8:eq:identidad-entrelazamiento-app-witt-u031}
\end{equation}
The direct sum defines a \(C_3\)-equivariant isomorphism
\begin{equation}
 T_\lambda:
 W_{\mathrm{APP}}\otimes\mathbb C
 \xrightarrow{\;\cong\;}
 W_c:=\bigoplus_{b=0}^{11}
 (A_{2,b}\otimes\mathbb C).
 \label{p12-83-c20-s8:eq:isomorfismo-app-witt-u031}
\end{equation}
\end{theorem}

\begin{proof}
For \(\mu=0\), the definition of \(g_b\) directly gives conjugation
by \(C\). For \(\mu=1\), one inserts
\(\mathsf RC\mathsf R=C^{-1}\). The
\cref{p12-83-c20-s8:lem:biyectividad-alineamiento-u031} guarantees that the direct sum
uses every Witt position exactly once and preserves the pair, sheet, and orbit.
\end{proof}

The theorem is relative to the simultaneous transport of the order of pairs, leaves, and
orbits, of the Witt origin, of the frames, and of the orientation. Changing a
single one of these coordinates does not define the same functor.

\subsection{Graded trace of the ternary generator}

Let
\[
 \rho_W:=\bigoplus_{b=0}^{11}g_b
\]
and consider the bosonic envelope
\begin{equation}
 \mathcal H_{\rm bos}(W_c)
 :=
 \operatorname{Sym}\!\left(
 \bigoplus_{n\ge1}W_c[n]
 \right).
 \label{p12-83-c20-s8:eq:fock-bosonico-doce-fibras-u031}
\end{equation}
If \(N\) is the degree operator, we define
\begin{equation}
 Z_{\rm Fock}(q)
 :=
 q^{-1}\operatorname{Tr}_{\mathcal H_{\rm bos}(W_c)}
 \left(\Gamma_{\rm bos}(\rho_W)q^N\right).
 \label{p12-83-c20-s8:eq:traza-graduada-fock-u031}
\end{equation}

\begin{proposition}[Cyclotomic polynomial of level three]
\label{p12-83-c20-s8:prop:producto-tres-doce-fibras-u031}
As a formal Laurent series,
\begin{equation}
 \boxed{
 Z_{\rm Fock}(q)
 =q^{-1}\prod_{n\ge1}(1+q^n+q^{2n})^{-12}
 =:t_3(q).}
 \label{p12-83-c20-s8:eq:t3-producto-u031}
\end{equation}
Its expansion begins with
\begin{equation}
 t_3(q)
 =q^{-1}-12+54q-76q^2-243q^3+1188q^4+O(q^5).
 \label{p12-83-c20-s8:eq:t3-expansion-u031}
\end{equation}
\end{proposition}

\begin{proof}
Each \(g_b\) is conjugate to \(C\), whence
\[
 \det(I-g_bq^n)=1+q^n+q^{2n}.
\]
The identity
\[
 \sum_{k\ge0}
 \operatorname{Tr}(\operatorname{Sym}^k\rho_W)t^k
 =\det(I-\rho_Wt)^{-1}
\]
applied in each degree produces the product. Formal multiplication through
degree five gives the expansion. In particular, the coefficient of \(q\) is
obtained as
\[
 [q^2]\,
 (1+q+q^2)^{-12}(1+q^2+q^4)^{-12}
 =66-12=54.
\]
\end{proof}

The value \(54\) appears here as the coefficient of a trace constructed from
twelve fibers. Unit U030 also obtained it as the oriented index of the
APP aggregate and proved that twelve is the unique positive solution of
cyclotomic--radial rigidity. The equality of the two values is a compatibility
between two realizations; neither uses the coefficient of \(J\) as input.

\subsection{Full-support word and fixed-point-free isometry}

Let
\[
 q_*=(1,1,1,1,1,1)\in\mathbb F_3^6.
\]
Multiplication by the Paley--Witt matrix gives
\begin{equation}
 q_*A_W=(2,1,1,1,1,1).
 \label{p12-83-c20-s8:eq:producto-qstar-aw-u031}
\end{equation}
Therefore,
\begin{equation}
 c_*:=(q_*,q_*A_W)
 =(1,1,1,1,1,1,2,1,1,1,1,1)
 \in C_W
 \label{p12-83-c20-s8:eq:palabra-soporte-total-u031}
\end{equation}
and all its coordinates are non-zero.

For each position \(b\), we fix
\begin{equation}
 P_b(c_*)=
 \begin{cases}
  \mathsf R,&(c_*)_b=1,\\
  I,&(c_*)_b=2,
 \end{cases}
 \qquad
 g_c:=\bigoplus_{b=0}^{11}P_b(c_*)CP_b(c_*)^{-1}.
 \label{p12-83-c20-s8:eq:isometria-orden-tres-u031}
\end{equation}

\begin{proposition}[Order, absence of fixed points and gluing]
\label{p12-83-c20-s8:prop:isometria-fijo-libre-pegado-u031}
The operator \(g_c\) has order three, has no nonzero fixed vectors, and
preserves the Niemeier lattice \(N(A_2^{12})\) constructed in U021.
\end{proposition}

\begin{proof}
Each block is conjugate to the Coxeter \(C\), whose minimal polynomial is
\(X^2+X+1\). Hence it has order three and has no eigenvalue one. The direct
sum preserves both properties. The Coxeter acts trivially on
\(A_2^*/A_2\); therefore each block preserves the discriminant class and the
union of lateral classes that defines the gluing.
\end{proof}

\subsection{Preservation of the Leech Neighbor}

Write
\[
 v:=v_{\mathcal F_*}
 =4\rho_1+\rho_2+\cdots+\rho_{12},
 \qquad
 v^2=54,
\]
for the vector selected by the incidence origin of
\eqref{p12-83-c20-s5:eq:vector-marco-incidencial}. Let
\[
 N:=N(A_2^{12}),
 \qquad
 N_0:=\ker(\langle\,\cdot\,,v\rangle\bmod3),
 \qquad
 N_v:=N_0+\mathbb Z\frac v3\cong\Lambda_{24}.
\]

\begin{theorem}[Neighbor-Oriented Preservation]
\label{p12-83-c20-s8:thm:preservacion-vecino-orden-tres-u031}
The displacements
\begin{equation}
 y_*:=\frac{g_cv-v}{3},
 \qquad
 z_*:=\frac{g_c^{-1}v-v}{3}
 \label{p12-83-c20-s8:eq:desplazamientos-vecino-u031}
\end{equation}
belong to \(N_0\). Consequently,
\[
 g_cN_0=N_0,
 \qquad
 g_cN_v=N_v.
\]
\end{theorem}

\begin{proof}
The discriminant words of \(y_*\) and \(z_*\) are
\(c_*\) and \(-c_*\), both in \(C_W\); hence
\(y_*,z_*\in N\). Computation in each factor \(A_2\) gives
\begin{equation}
 \langle y_*,v\rangle
 =\langle z_*,v\rangle
 =-(4^2+11)=-27\equiv0\pmod3.
 \label{p12-83-c20-s8:eq:emparejamientos-desplazamientos-u031}
\end{equation}
Thus both lie in \(N_0\). For \(x\in N_0\),
\[
 \langle g_cx,v\rangle
 =\langle x,g_c^{-1}v\rangle
 =\langle x,v\rangle+3\langle x,z_*\rangle
 \equiv0\pmod3.
\]
The same count with \(g_c^{-1}\) gives equality of the kernel. Finally,
\[
 g_c(v/3)=v/3+y_*,
\]
so that the class generated by the neighbor is preserved.
\end{proof}

\subsection{\(24\) orientations and lattice naturality}

The weight enumerator of \eqref{p12-83-c20-s1:eq:enumerador-pesos-cw} contains
\(24y^{12}\). Therefore, the set
\[
 C_W^\times:=
 \{u\in C_W:u_b\in\{1,2\}\text{ for every }b\}
\]
has 24 elements. For each \(u\in C_W^\times\), we define
\[
 P_b(u)=
 \begin{cases}
  \mathsf R,&u_b=1,\\
  I,&u_b=2,
 \end{cases}
 \qquad
 g_u:=\bigoplus_{b=0}^{11}P_b(u)CP_b(u)^{-1}.
\]

\begin{proposition}[Robustness of Full Support Orientations]
\label{p12-83-c20-s8:prop:robustez-orientaciones-u031}
For every \(u\in C_W^\times\),
\[
 \frac{g_uv-v}{3},
 \quad
 \frac{g_u^{-1}v-v}{3}
 \in N_0.
\]
In particular, each \(g_u\) preserves the neighbor constructed with the
transported orientation.
\end{proposition}

\begin{proof}
In simple root coordinates,
\[
 C\rho-\rho=(-2,-1),
 \qquad
 C^{-1}\rho-\rho=(-1,-2).
\]
The discriminant classes of the two displacements are \(u\) and \(-u\).
For the radial coefficient \(a_b\in\{4,1\}\),
\[
 \left\langle
 \frac{(C^{\pm1}-I)a_b\rho}{3},a_b\rho
 \right\rangle=-a_b^2.
\]
The sum over the twelve factors again gives \(-27\). The proof of
\cref{p12-83-c20-s8:thm:preservacion-vecino-orden-tres-u031} is applied component by
component.
\end{proof}

To type the recalibration, let
\(h\in\operatorname{Aut}(C_W)\) be monomial: a permutation \(\sigma\) of
the twelve coordinates followed by multipliers
\(\epsilon_b\in\{1,2\}\). We define its lattice lift
\begin{equation}
 \widetilde h
 :=P_\sigma\circ
 \bigoplus_{b=0}^{11}\mathsf R^{\nu_b},
 \qquad
 \nu_b=
 \begin{cases}
 0,&\epsilon_b=1,\\
 1,&\epsilon_b=2.
 \end{cases}
 \label{p12-83-c20-s8:eq:elevacion-monomial-u031}
\end{equation}
Its action on the discriminant is exactly \(h\). If the change simultaneously
transports the code, flag, incidence origin, vector \(v\), neighbor, and
frames, then
\begin{equation}
 g_{hu}=\widetilde h\,g_u\,\widetilde h^{-1}.
 \label{p12-83-c20-s8:eq:naturalidad-isometrias-u031}
\end{equation}
The ternary matrix \(h\) does not act without a lift on \(A_2^{12}\) or on
the Leech lattice.

\subsection{Lattice trace and computation of type \(0\)}

Let \(\widehat g_c\) be the standard lift of \(g_c\) to the lattice VOA
\[
 V_{N_v}=M(1)\otimes\mathbb C_\varepsilon[N_v].
\]
For \(j=1,2\), the oscillator--lattice decomposition produces
\begin{equation}
 \begin{split}
 Z_{0,j}(\tau)
 &:=
 \operatorname{Tr}_{V_{N_v}}
 \left(\widehat g_c^{\,j}q^{L_0-1}\right)\\
 &=q^{-1}
 \left(
 \sum_{\lambda\in N_v^{g_c^j}}
 \varepsilon_j(\lambda)
 q^{\langle\lambda,\lambda\rangle/2}
 \right)
 \prod_{n\ge1}\det(I-g_c^jq^n)^{-1}.
 \end{split}
 \label{p12-83-c20-s8:eq:factorizacion-traza-reticular-u031}
\end{equation}

\begin{proposition}[Trace of the nontrivial sectors]
\label{p12-83-c20-s8:prop:traza-reticular-t3-u031}
We have
\[
 Z_{0,1}(\tau)=Z_{0,2}(\tau)=t_3(\tau).
\]
\end{proposition}

\begin{proof}
The absence of fixed points implies
\[
 N_v^{g_c}=N_v^{g_c^2}=\{0\},
\]
so that the lattice sum in
\eqref{p12-83-c20-s8:eq:factorizacion-traza-reticular-u031} equals one. In each factor,
\(g_c^j\) has eigenvalues \(\zeta_3,\zeta_3^2\); therefore,
\[
 \det(I-g_c^jq^n)=1+q^n+q^{2n}.
\]
The twelve factors reproduce \eqref{p12-83-c20-s8:eq:t3-producto-u031}.
\end{proof}

\begin{theorem}[Twisted weight and type of the lift]
\label{p12-83-c20-s8:thm:peso-torcido-tipo-cero-u031}
The standard lift has order three, twisted weight
\begin{equation}
 \rho_{g_c}=\frac43
 \label{p12-83-c20-s8:eq:peso-torcido-u031}
\end{equation}
and type zero.
\end{theorem}

\begin{proof}
In
\(\mathfrak h=N_v\otimes_{\mathbb Z}\mathbb C\), each of the twelve
factors contributes one eigenvector with eigenvalue \(\zeta_3\) and another with eigenvalue
\(\zeta_3^2\). Therefore,
\[
 \dim\mathfrak h_{\zeta_3}
 =\dim\mathfrak h_{\zeta_3^2}=12.
\]
The formula for the weight of the twisted vacuum of order three gives
\[
 \rho_{g_c}
 =\frac1{4\cdot3^2}
 \bigl(1\cdot2\cdot12+2\cdot1\cdot12\bigr)
 =\frac43.
\]
For a lattice isometry of odd order, the standard elevation preserves
the order. The type is the class of
\[
 3^2\rho_{g_c}=12\equiv0\pmod3.
\]
Thus, the two initially independent hypotheses are calculated here.
\end{proof}

\subsection{Cyclic orbifold and dual class \texorpdfstring{\(3B\)}{3B}}

\begin{theorem}[Triadic orbifold of the Leech neighbor]
\label{p12-83-c20-s8:thm:orbifold-triadico-vnatural-u031}
The classical cyclic-orbifold theorems, applied to the standard lift
in \cref{p12-83-c20-s8:thm:peso-torcido-tipo-cero-u031}, yield
\begin{equation}
 \operatorname{Orb}_{\mathbb Z/3}
 (V_{N_v},\widehat g_c)
 \cong V^\natural.
 \label{p12-83-c20-s8:eq:orbifold-triadico-vnatural-u031}
\end{equation}
The dual automorphism belongs to the class \(3B\) and its series is
\begin{equation}
 T_{3B}(\tau)=t_3(\tau)+12.
 \label{p12-83-c20-s8:eq:serie-3b-u031}
\end{equation}
\end{theorem}

\begin{proof}
The isometry is fixed-point-free of order three; the twisted weight and type zero were calculated, not assumed. The cyclic-orbifold theorems identify the holomorphic extension with the Moonshine module and classify the dual automorphism as \(3B\) \cite{ChenLamShimakura2016,AbeLamYamada2017,Carnahan2017}. The identification uses these theorems after constructing the lattice object; it is not deduced solely from an equality of series.
\end{proof}

If \(V^{(i)}\) is the \(\widehat g_c^{\,i}\)-twisted module and
\[
 Z_{i,j}(\tau):=
 \operatorname{Tr}_{V^{(i)}}\!
 \left(
 \phi_i(\widehat g_c)^j q^{L_0-1}
 \right),
\]
the character of the orbifold is the average of the nine sectors:
\begin{equation}
 \operatorname{ch}_{V^\natural}(\tau)
 =\frac13\sum_{i,j\in\mathbb Z/3}Z_{i,j}(\tau)
 =J(\tau).
 \label{p12-83-c20-s8:eq:promedio-nueve-sectores-u031}
\end{equation}
The level three rational identity is
\begin{equation}
 J
 =t_3+12+\frac{10\,3^9}{t_3}
 +\frac{4\,3^{14}}{t_3^2}
 +\frac{3^{18}}{t_3^3}.
 \label{p12-83-c20-s8:eq:j-desde-t3-u031}
\end{equation}
Since \(t_3^{-1}=q+O(q^2)\),
\begin{equation}
 [q]J=54+10\cdot3^9=196884
 =54(5\cdot729+1).
 \label{p12-83-c20-s8:eq:coeficiente-196884-u031}
\end{equation}
This is a graded compatibility within the constructed corridor, not a
scalar morphism \(54\to V^\natural\).

\subsection{Lattice construction of order \(2\)}

The preceding route of order three and the classical construction of
Frenkel--Lepowsky--Meurman produce two presentations of the same moonshine VOA.
To declare the second, we fix the central extension
\begin{equation}
 1\longrightarrow\{\pm1\}
 \longrightarrow\widehat\Lambda_{24}
 \longrightarrow\Lambda_{24}
 \longrightarrow1
 \label{p12-83-c20-s8:eq:extension-central-leech-u031}
\end{equation}
with commutator
\((-1)^{\langle\lambda,\mu\rangle}\), and a cocycle
\(\varepsilon(\lambda,\mu)\) that represents it. The twisted group algebra is \(\mathbb C_\varepsilon[\Lambda_{24}]\), and
\begin{equation}
 V_\Lambda
 :=M(1)\otimes\mathbb C_\varepsilon[\Lambda_{24}]
 \label{p12-83-c20-s8:eq:voa-reticular-leech-u031}
\end{equation}
is the lattice VOA of central charge 24.

The negation \(\lambda\mapsto-\lambda\) has a lift
\(\theta\) compatible with \eqref{p12-83-c20-s8:eq:extension-central-leech-u031}. Its trace
is
\begin{equation}
 t_2(\tau)
 :=
 \operatorname{Tr}_{V_\Lambda}
 (\theta q^{L_0-1})
 =q^{-1}\prod_{n\ge1}(1+q^n)^{-24}.
 \label{p12-83-c20-s8:eq:t2-leech-u031}
\end{equation}
Indeed, negation pairs \(e^\lambda\) with \(e^{-\lambda}\), and the
absence of torsion leaves only \(\lambda=0\) in the lattice trace; each
of the 24 oscillators contributes \((1+q^n)^{-1}\).

Let \(V_\Lambda^T\) the modulo irreducible torcido by \(\theta\), and
\((V_\Lambda^T)^+\) its sector of signo positive. The construction FLM is
\begin{equation}
 V^\natural
 =(V_\Lambda)^+\oplus(V_\Lambda^T)^+.
 \label{p12-83-c20-s8:eq:vnatural-flm-u031}
\end{equation}
The route \(\mathbb Z/2\) of
\eqref{p12-83-c20-s8:eq:vnatural-flm-u031} and the route \(\mathbb Z/3\) of
\eqref{p12-83-c20-s8:eq:orbifold-triadico-vnatural-u031} are parallel constructions
identified by classical theorems; neither replaces the other.

\subsection{Automorphism group, character, and graded series}

eleven constructed \(V^\natural\), appear two publications distinct:
\begin{align}
 \operatorname{Aut}(V^\natural)&\cong\mathbb M,
 \label{p12-83-c20-s8:eq:aut-vnatural-monstruo-u031}\\
 \operatorname{ch}_{V^\natural}(\tau)
 &=J(\tau)=j(\tau)-744
 =q^{-1}+196884q+21493760q^2+\cdots.
 \label{p12-83-c20-s8:eq:caracter-vnatural-j-u031}
\end{align}
The former identifies the symmetry group; the latter counts graded
dimensions. The vacuum line is
\[
 (V^\natural)_0=\mathbb C\mathbf1,
\]
and the absence of roots in \(\Lambda_{24}\) gives
\[
 (V^\natural)_1=0.
\]
In degree two,
\[
 196884=1+196883=196560+324=54(5\cdot729+1).
\]
Only the first equality is a decomposition into dimensions of
representations of the Monster; the others are arithmetic controls.

For \(g\in\mathbb M\), the McKay--Thompson series is defined
\begin{equation}
 T_g(\tau):=
 \operatorname{Tr}_{V^\natural}
 \left(gq^{L_0-1}\right).
 \label{p12-83-c20-s8:eq:serie-mckay-thompson-u031}
\end{equation}

\begin{theorem}[Monstrous Moonshine]
\label{p12-83-c20-s8:thm:moonshine-monstruoso-u031}
The graded action of \(\mathbb M\) on \(V^\natural\) produces the family
\(\{T_g\}_{g\in\mathbb M}\). The Moonshine theorems identify these
series with genus-zero modular functions and prove their replicability
identities
\cite{ConwayNorton1979,Borcherds1992,FLM1988}.
\end{theorem}

This theorem has as input the pair
\[
 \bigl(V^\natural,\operatorname{Aut}(V^\natural)\bigr)
\]
and its grading. It is not the consequence of an arrow whose domain is the
abstract group \(\mathbb M\).

For \(J=T_{1A}\), let \(P_n\) be the Faber polynomials normalized by
\[
 P_n(J)=q^{-n}+O(q).
\]
The Hecke--Faber identity is
\begin{equation}
 P_n(J)=n\,T_nJ.
 \label{p12-83-c20-s8:eq:hecke-faber-u031}
\end{equation}
It controls all degrees; the coefficient 196884 is only the first
nontrivial manifestation.

\subsection{Paley–Witt–Golay–Leech–Moonshine Runner Dependency Graph}

\begin{figure}[htbp]
\centering
\begin{tikzcd}[column sep=huge,row sep=large]
 &C_W
   \arrow[dl,"W_{12}"']
   \arrow[d,"N(A_2^{12})"',font=\scriptsize]
   \arrow[dr,phantom,"\substack{\text{distinct}\\\text{branches}}",
          description,pos=.72,font=\scriptsize]&\\
 M_{12}
 &\Lambda_{24}
   \arrow[d,"V_\Lambda"]
 &(\mathcal G_{24},D,\ell)
   \arrow[d,"W_{24}"]\\
 &V^\natural
   \arrow[dl,"\operatorname{Aut}"']
   \arrow[d,"\operatorname{ch}"]
   \arrow[dr,"g\mapsto T_g"]
 &M_{24}\\
 \mathbb M
 &J
 &\{T_g\}_{g\in\mathbb M}
\end{tikzcd}
\caption{Exceptional bifurcation and three publications of
\(V^\natural\). The binary input \((\mathcal G_{24},D,\ell)\) is
independent; there is no arrow \(W_{24}\to\Lambda_{24}\).}
\label{p12-83-c20-s8:fig:grafo-excepcional-moonshine-u031}
\end{figure}

\begin{criterion}[Refutation criteria for the moonshine corridor]
\label{p12-83-c20-s8:crit:refutacion-corredor-lunar-u031}
The construction is refuted if any of the following facts occurs:
\begin{enumerate}
 \item \(\lambda\) is not bijective or
       \eqref{p12-83-c20-s8:eq:identidad-entrelazamiento-app-witt-u031} falla;
 \item \(c_*\notin C_W\), \(g_c\) has fixed points or does not preserve the
       neighbor;
 \item the displacements of
       \eqref{p12-83-c20-s8:eq:desplazamientos-vecino-u031} do not belong to \(N_0\);
 \item the lattice lift \(\widetilde h\) does not realize the monomial
       automorphism \(h\) on the discriminant;
 \item the twisted weight is not \(4/3\) or the lift is not of type zero;
 \item the orbifold is identified with \(V^\natural\) without verifying the
       hypotheses of the classical theorem;
 \item \(W_{24}\) is used as a constructor of the Leech lattice;
 \item an equality of dimensions is presented as a construction of
       \(\mathbb M\);
 \item the publications
       \(\operatorname{Aut}\), \(\operatorname{ch}\), and \(g\mapsto T_g\) are improperly serialized.
\end{enumerate}
\end{criterion}
